\documentclass{article}

% ============================================================
% TEMPLATE DO PROJETO DESENVOLVIDO (ETAPA 2)
% Adaptado do template oficial do NeurIPS 2026 para uso na disciplina.
%
% INSTRUÇÕES PARA OS ALUNOS:
% - Baixe também o arquivo neurips_2026.sty do template original em:
%   https://www.overleaf.com/latex/templates/formatting-instructions-for-neurips-2026/bjdwqfdkyftc
% - LIMITE DE PÁGINAS DESTA ETAPA: tentem manter entre 4 a 6 páginas, sem contar as referências, figuras e tabelas. 
% - Este documento parte da proposta (Etapa 1): reaproveitem e atualizem
%   as seções Introduction, Related Work, Data e Methodology conforme o
%   que de fato foi feito, e acrescentem os resultados obtidos.
% - Não é necessário anonimizar o documento nem incluir checklist.
% ============================================================

\usepackage[preprint]{neurips_2026}

\usepackage[utf8]{inputenc}
\usepackage[T1]{fontenc}
\usepackage{hyperref}
\usepackage{url}
\usepackage{booktabs}
\usepackage{amsfonts}
\usepackage{nicefrac}
\usepackage{microtype}
\usepackage{xcolor}
\usepackage{graphicx}

\title{Título do Projeto: [substituam pelo título do projeto]}

\author{%
  Nome do(a) Aluno(a) 1 \\
  Código do Aluno(a) 1 \\
  \texttt{aluno1@email.com} \\
  \And
  Nome do(a) Aluno(a) 2 \\
  Código do Aluno(a) 2 \\
  \texttt{aluno1@email.com} \\
}
\begin{document}

\maketitle

\begin{abstract}
  Resumo de 4 a 6 linhas: problema, método usado, principais resultados
  (com números, se possível) e conclusão principal. Diferente do resumo
  da proposta, este já deve mencionar resultados obtidos. \\
  O documento pode ser escrito em português ou em inglês — escolham um e mantenham a consistência até o fim. Termos técnicos consagrados na literatura (self-attention, embedding, fine-tuning, dataset, etc.) podem permanecer em inglês mesmo em textos escritos em português; não inventem traduções pouco usuais para eles (evitem, por exemplo, "autoatenção" ou "ajuste fino" quando o termo em inglês é o que todo mundo usa e entende).
\end{abstract}

\section{Introduction}
\label{sec:intro}
% Atualizem a introdução da proposta, se necessário. Mantenham a
% pergunta de pesquisa e os objetivos, ajustando caso o escopo do
% projeto tenha mudado ao longo do desenvolvimento.

\section{Related Work}
\label{sec:related}
% Atualizem a revisão da proposta com trabalhos adicionais encontrados
% ao longo do desenvolvimento, se houver.

\section{Data}
\label{sec:data}
% Descrevam os dados efetivamente usados (podem diferir do que foi
% planejado na proposta) e o pré-processamento efetivamente aplicado.
% Se algo mudou em relação à proposta, expliquem por quê.

\section{Methodology}
\label{sec:method}
% Descrevam o método/modelo efetivamente implementado. Se ele mudou em
% relação à proposta, expliquem as razões da mudança.

\section{Experiments}
\label{sec:exp}
% O que escrever aqui:
% - Desenho experimental efetivamente executado.
% - Baselines usados na comparação.
% - Métricas de avaliação usadas.
% - Detalhes de implementação relevantes para reprodutibilidade
%   (hiperparâmetros, hardware, tempo de treinamento, splits de dados).

\section{Results}
\label{sec:results}
% O que escrever aqui:
% - Apresentem os resultados obtidos de forma objetiva (tabelas, figuras,
%   gráficos). Ver exemplos de tabela e figura comentados abaixo.
% - Não interpretem os resultados aqui — isso vai na seção Discussion.

% \begin{table}[h]
%   \caption{Resultados obtidos no conjunto de teste.}
%   \label{tab:results}
%   \centering
%   \begin{tabular}{lcc}
%     \toprule
%     Método & Métrica 1 & Métrica 2 \\
%     \midrule
%     Baseline & & \\
%     Proposto & & \\
%     \bottomrule
%   \end{tabular}
% \end{table}

% \begin{figure}[h]
%   \centering
%   \includegraphics[width=0.7\linewidth]{resultados.pdf}
%   \caption{Descrição do que a figura mostra e principal conclusão.}
%   \label{fig:results}
% \end{figure}

\section{Discussion}
\label{sec:discussion}
% O que escrever aqui:
% - Interpretem os resultados da Seção \ref{sec:results}: o que eles
%   significam em relação à pergunta de pesquisa?
% - Comparem com o baseline e com a literatura discutida em Related Work.
% - Casos de erro/falha interessantes, se houver.

\section{Limitations}
\label{sec:limitations}
% O que escrever aqui:
% - Limitações dos dados, do método ou da avaliação.
% - O que não foi possível testar por restrições de tempo/recursos.
% - Ameaças à validade dos resultados (ex.: tamanho da amostra,
%   generalização, vazamento de dados).

\section{Conclusion}
\label{sec:conclusion}
% O que escrever aqui:
% - Retomem a pergunta de pesquisa e respondam-na com base nos resultados.
% - Principais conclusões do trabalho, em 3 a 5 frases.
% - Trabalhos futuros: o que fariam com mais tempo/recursos.

\bibliographystyle{plain}
\bibliography{references}

\end{document}
